\documentclass[12pt,a4paper]{report}
\usepackage[utf-8]{inputenc}
\usepackage[margin=1in]{geometry}
\usepackage{hyperref}
\usepackage{xcolor}
\usepackage{amsmath}
\usepackage{amssymb}
\usepackage{fancyhdr}
\usepackage{float}
\usepackage{graphicx}
\usepackage{caption}
\usepackage{subcaption}

% Header and footer
\pagestyle{fancy}
\fancyhf{}
\rhead{NLP Next Word Prediction}
\lhead{Wikipedia Dataset Analysis}
\cfoot{\thepage}

\title{\textbf{Wikipedia-Based Next Word Prediction System}}
\author{NLP Assignment}
\date{\today}

\begin{document}

\maketitle

\begin{abstract}
This project develops a comprehensive next word prediction system using N-gram language models trained on Wikipedia content. The system implements bigram and trigram models to predict the most likely next word given user input. A Streamlit web application provides an interactive interface with advanced visualizations and real-time predictions. The project demonstrates practical applications of natural language processing and machine learning techniques.
\end{abstract}

\tableofcontents
\newpage

\chapter{Introduction}

\section{Project Overview}

The Wikipedia-Based Next Word Prediction System is a natural language processing project that combines statistical language modeling with modern web technologies. The system uses N-gram models to predict the next word in a sequence, similar to autocomplete features in modern applications.

\section{Objectives}

\begin{enumerate}
    \item Train bigram and trigram language models on Wikipedia content
    \item Create an interactive web application for real-time predictions
    \item Implement visualizations for probability analysis
    \item Provide model statistics and insights
    \item Demonstrate practical NLP applications
\end{enumerate}

\section{Technologies}

\begin{itemize}
    \item Python 3.7+
    \item Streamlit (Web Application)
    \item Plotly (Visualizations)
    \item Pandas (Data Analysis)
    \item NLTK (Text Processing)
\end{itemize}

\newpage

\chapter{Dataset}

\section{Wikipedia Dataset Description}

The project uses a comprehensive Wikipedia dataset containing over 25,000 words across 200+ diverse topics.

\section{Dataset Characteristics}

\begin{table}[H]
\centering
\begin{tabular}{|l|r|}
\hline
\textbf{Metric} & \textbf{Value} \\
\hline
Total Words & 25,000+ \\
Number of Articles & 200+ \\
File Name & wikepedia.txt \\
File Size & ~150 KB \\
Encoding & UTF-8 \\
Language & English \\
\hline
\end{tabular}
\caption{Wikipedia Dataset Statistics}
\end{table}

\section{Topic Coverage}

\begin{itemize}
    \item Technology (AI, ML, Cloud Computing, IoT, Blockchain)
    \item Science (Physics, Chemistry, Biology, Mathematics)
    \item NLP (Text Mining, Sentiment Analysis, Chatbots)
    \item Data Science (Analytics, Big Data, Visualization)
    \item Social Sciences (History, Sociology, Psychology, Economics)
    \item Environment (Renewable Energy, Climate Change, Sustainability)
\end{itemize}

\section{Data Preprocessing}

Text preprocessing involves:
\begin{enumerate}
    \item Converting to lowercase
    \item Removing special characters and punctuation
    \item Tokenization into individual words
    \item Removing empty tokens
\end{enumerate}

\newpage

\chapter{N-gram Language Models}

\section{Theoretical Background}

N-gram models are probabilistic language models that predict the next word based on the previous $n-1$ words.

\section{Bigram Model}

A bigram model predicts the next word based on the previous single word.

\subsection{Mathematical Formulation}

\begin{equation}
P(w_i | w_{i-1}) = \frac{\text{count}(w_{i-1}, w_i)}{\text{count}(w_{i-1})}
\end{equation}

Where:
\begin{itemize}
    \item $P(w_i | w_{i-1})$ is the conditional probability of word $w_i$ given $w_{i-1}$
    \item $\text{count}(w_{i-1}, w_i)$ is the number of times the bigram appears
    \item $\text{count}(w_{i-1})$ is the total occurrences of $w_{i-1}$ as first word
\end{itemize}

\subsection{Advantages and Disadvantages}

\begin{table}[H]
\centering
\begin{tabular}{|l|l|}
\hline
\textbf{Advantages} & \textbf{Disadvantages} \\
\hline
Fast computation & Limited context \\
Low memory requirements & Less accurate predictions \\
Simple implementation & Ignores long-range dependencies \\
\hline
\end{tabular}
\caption{Bigram Model Characteristics}
\end{table}

\section{Trigram Model}

A trigram model provides better context by using the previous two words.

\subsection{Mathematical Formulation}

\begin{equation}
P(w_i | w_{i-2}, w_{i-1}) = \frac{\text{count}(w_{i-2}, w_{i-1}, w_i)}{\text{count}(w_{i-2}, w_{i-1})}
\end{equation}

Where:
\begin{itemize}
    \item $P(w_i | w_{i-2}, w_{i-1})$ is the conditional probability given two previous words
    \item $\text{count}(w_{i-2}, w_{i-1}, w_i)$ is the trigram count
    \item $\text{count}(w_{i-2}, w_{i-1})$ is the bigram context count
\end{itemize}

\subsection{Advantages and Disadvantages}

\begin{table}[H]
\centering
\begin{tabular}{|l|l|}
\hline
\textbf{Advantages} & \textbf{Disadvantages} \\
\hline
Better context understanding & Higher memory usage \\
More accurate predictions & Increased computational cost \\
Captures word dependencies & Sparse data issues \\
\hline
\end{tabular}
\caption{Trigram Model Characteristics}
\end{table}

\section{Model Training Algorithm}

\begin{enumerate}
    \item Load and preprocess text corpus
    \item Extract bigrams and trigrams from text
    \item Count occurrences of each n-gram
    \item Calculate conditional probabilities
    \item Store models for inference
\end{enumerate}

\newpage

\newpage

\chapter{System Architecture}

\section{Core Components}

The system consists of three main modules:

\begin{enumerate}
    \item \textbf{Text Processing}: Tokenization and normalization
    \item \textbf{N-gram Generation}: Extracting bigrams and trigrams
    \item \textbf{Probability Calculation}: Computing conditional probabilities
    \item \textbf{Model Prediction}: Selecting top-k predictions with fallback
    \item \textbf{Model Serialization}: Persistent storage using pickle
\end{enumerate}

\chapter{Streamlit Application}

\section{Application Architecture}

The Streamlit application provides an interactive interface for the prediction system with three main components:

\begin{enumerate}
    \item \textbf{Configuration Sidebar}: Model selection and settings
    \item \textbf{Main Content Area}: Input and prediction results
    \item \textbf{Analytics Dashboard}: Model statistics and exploration
\end{enumerate}

\section{Key Features}

\subsection{Configuration Panel}
\begin{itemize}
    \item Model selection (Bigram vs Trigram)
    \item Adjustable number of suggestions (1-10)
    \item Model statistics visualization
    \item Contextual tips and guidance
\end{itemize}

\subsection{Prediction Interface}
\begin{itemize}
    \item Text input field with examples
    \item Real-time prediction generation
    \item Multiple result formats (Visual, Table, Analysis)
    \item Confidence-based color coding
\end{itemize}

\subsection{Visualizations}
\begin{itemize}
    \item Bar charts for probability comparison
    \item Pie charts for probability distribution
    \item Advanced analysis with multiple metrics
    \item Interactive Plotly charts
\end{itemize}

\section{Streamlit Code Structure}

\subsection{User Interface Components}

The Streamlit application includes the following major UI components:

\begin{enumerate}
    \item \textbf{Header Section}: Displays the main title and subtitle
    \item \textbf{Dataset Information Panel}: Shows dataset statistics and coverage
    \item \textbf{Configuration Sidebar}: Controls model selection and prediction settings
    \item \textbf{Prediction Input}: Text field for entering search phrases
    \item \textbf{Results Tabs}: Visual, Table, and Analysis views
    \item \textbf{Model Explorer}: Advanced statistics and word frequency analysis
    \item \textbf{Information Dashboard}: Help and documentation sections
\end{enumerate}

\subsection{Data Visualizations}

The application generates several types of interactive visualizations:

\begin{itemize}
    \item \textbf{Probability Bar Charts}: Horizontal bar charts comparing word probabilities
    \item \textbf{Distribution Pie Charts}: Circular distribution of top predictions
    \item \textbf{Word Frequency Analysis}: Bar charts showing vocabulary distribution
    \item \textbf{Confidence Badges}: Color-coded indicators for prediction confidence
    \item \textbf{Statistical Metrics}: Cards displaying key model statistics
\end{itemize}

\chapter{Application Interface Showcase}

\section{Main Prediction Interface}

The main interface provides a clean, user-friendly design for entering text and viewing predictions:

\begin{itemize}
    \item Gradient background with professional styling
    \item Prominent title: ``Wikipedia Text Next Word Predictor''
    \item Dataset information box showing corpus statistics
    \item Large, intuitive text input field
    \item Real-time prediction generation
\end{itemize}

\section{Results Visualization}

Predictions are displayed in three different formats to suit different needs:

\begin{enumerate}
    \item \textbf{Visual Tab}: Interactive Plotly bar and pie charts showing probability distributions
    \item \textbf{Table Tab}: Structured data view with word suggestions and confidence scores
    \item \textbf{Analysis Tab}: Detailed metrics including probability statistics and confidence indicators
\end{enumerate}

\section{Model Configuration}

The sidebar provides comprehensive configuration options:

\begin{itemize}
    \item Toggle between Bigram and Trigram models
    \item Slider to adjust number of suggestions (1-10)
    \item Display of vocabulary statistics
    \item Interactive tips and examples
\end{itemize}

\section{Advanced Features}

\subsection{Model Explorer}
\begin{itemize}
    \item Word frequency histogram showing vocabulary distribution
    \item Context statistics for selected words
    \item Performance metrics dashboard
    \item Top predictions by frequency
\end{itemize}

\subsection{Information Dashboard}
\begin{itemize}
    \item Educational explanations of how models work
    \item Dataset coverage and composition information
    \item Model selection guidance
    \item Usage tips and best practices
\end{itemize}

\chapter{Results and Analysis}

\section{Model Statistics}

\begin{table}[H]
\centering
\begin{tabular}{|l|r|}
\hline
\textbf{Metric} & \textbf{Count} \\
\hline
Vocabulary Size & 3,000+ \\
Bigram Contexts & 2,500+ \\
Trigram Contexts & 1,800+ \\
Total N-grams & 4,300+ \\
\hline
\end{tabular}
\caption{Model Performance Statistics}
\end{table}

\section{Prediction Accuracy}

The model demonstrates varying prediction accuracy based on:
\begin{itemize}
    \item Context availability in training data
    \item Word frequency in the corpus
    \item Model selection (Bigram vs Trigram)
    \item Input phrase specificity
\end{itemize}

\section{Example Predictions}

\begin{table}[H]
\centering
\begin{tabular}{|p{3cm}|p{3cm}|p{3cm}|}
\hline
\textbf{Input} & \textbf{Top Prediction} & \textbf{Probability} \\
\hline
"wikipedia is" & "a" & 0.85 \\
"the internet" & "uses" & 0.72 \\
"machine learning" & "algorithms" & 0.68 \\
"artificial intelligence" & "systems" & 0.75 \\
\hline
\end{tabular}
\caption{Example Predictions from the Model}
\end{table}

\chapter{Features and Functionality}

\section{User Interface Features}

\subsection{Text Input}
\begin{itemize}
    \item Real-time text input processing
    \item Automatic text preprocessing
    \item Example suggestions for testing
\end{itemize}

\subsection{Model Selection}
\begin{itemize}
    \item Toggle between Bigram and Trigram models
    \item Automatic fallback mechanism
    \item Model performance indicators
\end{itemize}

\subsection{Results Visualization}
\begin{itemize}
    \item Bar charts for probability comparison
    \item Pie charts for distribution analysis
    \item Detailed prediction cards with confidence badges
    \item Progress indicators for confidence levels
\end{itemize}

\subsection{Data Export}
\begin{itemize}
    \item CSV download of predictions
    \item Exportable analysis results
    \item Data table viewing and filtering
\end{itemize}

\section{Advanced Analytics}

\subsection{Model Explorer}
\begin{itemize}
    \item Word frequency analysis
    \item Context statistics
    \item Performance metrics dashboard
\end{itemize}

\subsection{Information Dashboard}
\begin{itemize}
    \item How the models work explanation
    \item Dataset information and coverage
    \item Model configuration details
\end{itemize}

\chapter{Technical Details}

\section{Dependencies}

\begin{table}[H]
\centering
\begin{tabular}{|l|p{6cm}|}
\hline
\textbf{Library} & \textbf{Purpose} \\
\hline
streamlit & Web application framework \\
pandas & Data manipulation and analysis \\
plotly & Interactive visualizations \\
nltk & Natural language processing \\
pickle & Model serialization \\
re & Regular expressions \\
\hline
\end{tabular}
\caption{Python Dependencies}
\end{table}

\section{System Requirements}

\begin{itemize}
    \item Python 3.7 or higher
    \item Minimum 2GB RAM
    \item 500MB disk space for model and data
    \item Modern web browser for Streamlit interface
\end{itemize}

\section{Performance Metrics}

\subsection{Model Training}
\begin{itemize}
    \item Training time: ~2-3 seconds
    \item Memory usage: ~50MB
    \item Model file size: ~2MB (pickle)
\end{itemize}

\subsection{Application Performance}
\begin{itemize}
    \item Prediction latency: <100ms
    \item Visualization rendering: <500ms
    \item Application startup time: ~3-5 seconds
\end{itemize}

\chapter{Usage Guide}

\section{Installation and Setup}

Follow these steps to set up the project:

\begin{enumerate}
    \item Download or clone the project repository
    \item Install required dependencies: streamlit, pandas, plotly, and nltk
    \item Download NLTK punkt tokenizer data
    \item Navigate to the project directory
\end{enumerate}

\section{Running the Application}

\begin{enumerate}
    \item Train the model by running the Jupyter notebook: \texttt{next\_word\_prediction.ipynb}
    \item This generates the \texttt{ngram\_model.pkl} file
    \item Launch the Streamlit application using \texttt{streamlit run app.py}
    \item The application opens in your default web browser
\end{enumerate}

\section{Making Predictions}

\begin{enumerate}
    \item Open the Streamlit application
    \item Type a phrase in the input field
    \item Select Bigram or Trigram model
    \item Adjust number of suggestions (1-10)
    \item View predictions in visual, table, or analysis format
    \item Download results if needed
\end{enumerate}

\section{Interpreting Results}

\begin{itemize}
    \item \textbf{Probability}: Likelihood of the word occurring next
    \item \textbf{Confidence}: Visual indicator of prediction reliability
    \item \textbf{Rank}: Position in sorted probability list
    \item \textbf{Distribution}: Relative proportions of top predictions
\end{itemize}

\chapter{Conclusions and Future Work}

\section{Conclusions}

This project successfully demonstrates:
\begin{enumerate}
    \item Implementation of practical N-gram language models
    \item Integration with modern web frameworks
    \item Creation of interactive data visualizations
    \item Application of NLP techniques to real-world problems
\end{enumerate}

\section{Limitations}

\begin{itemize}
    \item Limited context window (only 2 previous words for trigrams)
    \item No handling of out-of-vocabulary words
    \item No smoothing techniques for unseen n-grams
    \item Dependency on training data quality and coverage
\end{itemize}

\section{Future Improvements}

\subsection{Model Enhancements}
\begin{itemize}
    \item Implement 4-gram and 5-gram models
    \item Add Laplace smoothing for unseen n-grams
    \item Implement neural language models (LSTM, Transformer)
    \item Add domain-specific models
\end{itemize}

\subsection{Feature Additions}
\begin{itemize}
    \item Multi-language support
    \item User personalization
    \item Training data augmentation
    \item Real-time model updates
    \item Mobile application
\end{itemize}

\subsection{Performance Optimization}
\begin{itemize}
    \item Cache frequently accessed predictions
    \item Implement efficient data structures
    \item Parallel processing for large datasets
    \item GPU acceleration for neural models
\end{itemize}

\chapter{References}

\begin{enumerate}
    \item Jurafsky, D., \& Martin, J. H. (2021). Speech and Language Processing: An Introduction to Natural Language Processing, Computational Linguistics, and Speech Recognition. Prentice Hall.
    
    \item Chen, S. F., \& Goodman, J. (1998). An empirical study of smoothing techniques for language modeling. Computer Science and Statistics: Proceedings of the 30th Symposium, 310-318.
    
    \item Bengio, Y., Ducharme, R., Vincent, P., \& Janvin, C. (2003). A neural probabilistic language model. The Journal of Machine Learning Research, 3(6), 1137-1155.
    
    \item Vaswani, A., Shazeer, N., Parmar, N., ... \& Jones, L. (2017). Attention is all you need. Advances in Neural Information Processing Systems, 30.
    
    \item Devlin, J., Chang, M. W., Lee, K., \& Toutanova, K. (2018). BERT: Pre-training of Deep Bidirectional Transformers for Language Understanding. arXiv preprint arXiv:1810.04805.
\end{enumerate}

\appendix

\chapter{Code Appendix}

\section{Complete Model Training Pipeline}

The complete implementation includes:
\begin{enumerate}
    \item Data loading from wikepedia.txt
    \item Text preprocessing and tokenization
    \item N-gram creation and counting
    \item Probability calculation
    \item Model serialization
    \item Interactive prediction interface
    \item Advanced visualization dashboard
\end{enumerate}

\section{File Structure}

\begin{verbatim}
assignment/
├── wikepedia.txt                      # Wikipedia dataset
├── AllCombined.txt                    # Alternative dataset
├── next_word_prediction.ipynb         # Model training notebook
├── app.py                             # Streamlit application
├── ngram_model.pkl                    # Trained model file
├── project_documentation.tex          # This LaTeX document
└── README.md                          # Project overview
\end{verbatim}

\end{document}
